\documentclass[11pt,letterpaper]{article}
\usepackage[margin=1in]{geometry}
\usepackage{booktabs}
\usepackage{longtable}
\usepackage{hyperref}
\usepackage{xcolor}
\usepackage{enumitem}
\usepackage{tabularx}

\title{Replication Report (Re-pass, LaTeX render):\\
Shrestha et al.\ 2022 --- Trehalose Pathway Prediction in \textit{Variovorax} sp.\ PAMC28711}
\author{BVBRC Replication Project (BVBRC-04)}
\date{Re-pass 2026-06-23; LaTeX render 2026-07-05}

\begin{document}
\maketitle

\begin{abstract}
Independent re-pass replication of Shrestha et al.\ (\textit{BMC Genomic Data} 23:4, 2022),
which compares KEGG, MetaCyc, and RAST annotations of the trehalose-biosynthesis
enzyme repertoire in \textit{Variovorax} sp.\ PAMC28711 (NZ\_CP014517.1). Verdict:
\textbf{PARTIAL}. Of 37 enumerated paper claims, 26 are testable on free
compute; 21 verified, 4 partial, 0 contradicted. The paper's central methodological
thesis (annotation-system divergence at the TreY step) reproduces cleanly and is
independently corroborated by PGAP's pseudogene flag on the same locus. The
paper's broader biological framing ``three functional trehalose-biosynthesis
pathways'' is weakened by that same pseudogene call, which the paper does not
address. MetaCyc rows and 2018 database-snapshot statistics remain permanently
un-testable on free compute.
\end{abstract}

\section{Paper Reference}
\begin{itemize}[leftmargin=1.2em,itemsep=1pt]
  \item \textbf{Citation:} Shrestha P, Kim M-S, Elbasani E, Kim J-D, Oh T-J.
    \textit{BMC Genomic Data} 23:4 (2022).
  \item \textbf{DOI:} \href{https://doi.org/10.1186/s12863-021-01020-y}{10.1186/s12863-021-01020-y} (CC BY 4.0)
  \item \textbf{PMID / PMC:} 34991451 / PMC8734048
  \item \textbf{Genome:} NZ\_CP014517.1 / CP014517.1 (single circular chromosome, 4{,}316{,}152 bp)
  \item \textbf{KEGG organism code:} \texttt{vaa}
  \item \textbf{BV-BRC genome ID:} 1795631.3
  \item \textbf{NCBI Assembly:} GCA\_001577265.1 (RefSeq GCF\_001577265.1)
  \item \textbf{PDF SHA-256:} \texttt{f0f7a5addf671072cdab447b7c4b3b42c9bb07472e95305f9aba957a18ffc424}
    (\texttt{paper/shrestha2022.pdf}, 1.84~MB)
\end{itemize}

\section{Parser Provenance and Reproducer Pipeline}
\begin{tabularx}{\linewidth}{@{}l l X@{}}
\toprule
Step & Tool & Output \\
\midrule
PDF $\to$ text        & \texttt{pdftotext -layout} (poppler 25.x) & \texttt{paper/shrestha2022.txt} (361 lines) \\
GenBank $\to$ features & BioPython \texttt{SeqIO.parse('genbank')} & \texttt{results/repass/genome\_features.json} \\
KEGG REST cross-check  & \texttt{urllib} $\to$ \texttt{rest.kegg.jp} & \texttt{results/repass/kegg\_crosscheck.json} \\
Trehalose / glycogen scan & BioPython + product-name regex & \texttt{results/repass/trex\_and\_cluster.json} \\
BV-BRC metadata        & \texttt{urllib} $\to$ \texttt{bv-brc.org/api/genome} & \texttt{results/repass/bvbrc\_metadata.json} \\
MetaCyc                & \textbf{BLOCKED} --- no PAMC28711 PGDB & (unchanged from pass-1) \\
\bottomrule
\end{tabularx}

\medskip
All code under \texttt{code/repass/} is plain Python 3, uses only BioPython + \texttt{urllib}, runs
on a laptop in seconds, uses only free public APIs, and involves no LLM calls in the numeric path.
All inputs are hashed in \texttt{results/repass/parser\_provenance.json}. The cached
\texttt{data/CP014517.1.gb} carries zero \texttt{/EC\_number} qualifiers (PGAP keys on product
names), so re-pass scripts key on product-name regexes and validate against KEGG KO links.

\section{Scope of the Paper}
Shrestha 2022 is a short replication-style methods paper that:
\begin{enumerate}[leftmargin=1.5em,itemsep=1pt]
  \item Identifies trehalose-biosynthesis enzymes in \textit{Variovorax} sp.\ PAMC28711 via three annotation
        systems (KEGG \texttt{vaa}, MetaCyc, RAST/SEED Viewer).
  \item Presents a $5\times 3$ presence/absence comparison (Table 1) for OtsA-OtsB, TreY/TreZ, and TreS pathways.
  \item Reports 2018-snapshot database statistics (Table 2: 2{,}688 / 339 / 381 / 530 pathways/maps; 15{,}329 vs 11{,}004 reactions).
  \item Concludes KEGG and MetaCyc both miss TreY (EC 5.4.99.15) while RAST finds it.
\end{enumerate}
Testable units enumerated in re-pass: \textbf{37 claims}. Of those, 11 are inherently untestable on free compute
(5 MetaCyc rows $+$ 6 historical Table-2/Discussion stats), leaving \textbf{26 testable}.

\section{Methods (Re-pass Additions)}
\begin{tabularx}{\linewidth}{@{}l X X@{}}
\toprule
New in re-pass & What it does & Why pass-1 did not have it \\
\midrule
\texttt{01\_genome\_features.py} & Recomputes size, GC\%, CDS/rRNA/tRNA/pseudogene counts, isolation source directly from GenBank & Pass-1 quoted these from BV-BRC without recomputation \\
\texttt{02\_trex\_and\_full\_cluster.py} & Adds TreX (EC 3.2.1.68), dumps full trehalose + glycogen cluster regions & Fig 2A names TreX; pass-1 stopped at TreY/TreZ/TreS \\
\texttt{03\_kegg\_crosscheck.py} & Resolves all paper EC$\to$KO, queries \texttt{link/vaa/ko:K\ldots} per KO & Pass-1 used KEGG but did not record per-KO link queries reproducibly \\
\texttt{04\_bvbrc\_metadata.py} & Verifies isolation country=Antarctica, host=\textit{Himantormia} (lichen), assembly accession, genome status & Pass-1 took Antarctic/lichen claim on trust from local notes \\
\bottomrule
\end{tabularx}

\section{Results}

\subsection{Genome features (from PGAP GenBank, directly recomputed)}
\begin{tabularx}{\linewidth}{@{}l l X@{}}
\toprule
Metric & Re-pass value & Source \\
\midrule
Total length     & 4{,}316{,}152 bp             & \texttt{data/CP014517.1.gb} \\
Records / contigs& 1 (circular)                 & PGAP \\
GC\%             & 65.973                       & PGAP sequence \\
CDS (total)      & 4{,}104                      & PGAP \\
\quad of which \texttt{/pseudo} & 129           & PGAP \\
rRNA             & 6 (2$\times$5S+2$\times$16S+2$\times$23S) & PGAP \\
tRNA             & 46                           & PGAP \\
Host             & \textit{Himantormia} (lichen)& PGAP \texttt{/host} qualifier \\
Country          & Antarctica                   & BV-BRC \texttt{isolation\_country} \\
Collection date  & 2015                         & PGAP \texttt{/collection\_date} \\
Assembly         & GCA\_001577265.1             & BV-BRC \\
BioSample        & SAMN04457487                 & BV-BRC \\
BV-BRC PATRIC CDS& 4{,}263                      & BV-BRC API \\
\bottomrule
\end{tabularx}

\subsection{Trehalose-pathway enzyme presence --- re-pass}
\begin{tabularx}{\linewidth}{@{}l l l l l X l@{}}
\toprule
Enzyme & EC & Pap.KEGG & Pap.MC & Pap.RAST & Re-pass (KO $\to$ vaa / PGAP locus) & Status \\
\midrule
OtsA    & 2.4.1.15  & O & O & O & K00697 $\to$ AX767\_06265; +strand 1{,}238{,}237..1{,}239{,}625 & \checkmark \\
OtsB    & 3.1.3.12  & O & O & O & K01087 $\to$ AX767\_06260; +strand 1{,}237{,}485..1{,}238{,}237 & \checkmark \\
TreY    & 5.4.99.15 & X & X & O & K06044 $\to$ NONE; \textbf{PSEUDO} AX767\_16200 (frameshift) & \checkmark \\
TreZ    & 3.2.1.141 & O & O & O & K01236 $\to$ AX767\_16205; functional & \checkmark \\
TreS    & 5.4.99.16 & O & O & O & K05343 $\to$ AX767\_16215 (dual EC 5.4.99.16 / 3.2.1.1) & \checkmark \\
TreF/H  & 3.2.1.28  & --- & --- & O & K01194 $\to$ AX767\_10110 (\texttt{treF}) & \checkmark \\
TreP    & 2.4.1.64  & --- & --- & --- & K00691 $\to$ NONE (absent) & \checkmark \\
TreT    & 2.4.1.245 & --- & --- & --- & K13057 $\to$ NONE (absent) & \checkmark \\
TreX (new) & 3.2.1.68 & --- & (Fig 2A) & --- & K02438 $\to$ AX767\_11830; K01214 $\to$ AX767\_10865 (2 copies) & \checkmark \\
\bottomrule
\end{tabularx}

\subsection{Glycogen biosynthesis cluster ($\sim$2.41 Mbp) --- new in re-pass}
\begin{tabularx}{\linewidth}{@{}l X l l@{}}
\toprule
Locus & Product & KO & EC \\
\midrule
AX767\_11815 & glycogen phosphorylase                  & K00688 & 2.4.1.1 \\
AX767\_11820 & glycogen synthase                       & K00703 & 2.4.1.21 \\
AX767\_11825 & glucose-1-P adenylyltransferase         & K00975 & 2.7.7.27 \\
AX767\_11830 & glycogen debranching enzyme (TreX cand.)& K02438 & 3.2.1.68 \\
AX767\_11845 & glycogen-branching enzyme               & K00700 & 2.4.1.18 \\
\bottomrule
\end{tabularx}

\noindent The glycogen biosynthesis operon is complete and functional. MetaCyc trehalose
biosynthesis V (PWY-2661) consumes glycogen $\to$ maltodextrin $\to$ TreY substrate, so the
TreY/TreZ pathway has all upstream infrastructure except the TreY step itself.

\subsection{KEGG \texttt{vaa} coverage (new)}
\begin{itemize}[leftmargin=1.2em,itemsep=1pt]
  \item 4{,}159 entries in KEGG \texttt{vaa} (\texttt{list/vaa}).
  \item 3{,}975 proteins; 2{,}351 have KO assignments.
  \item 133 KEGG pathway maps with vaa genes (\texttt{list/pathway/vaa}).
  \item 56 KEGG modules with vaa genes (\texttt{link/module/vaa}, unique).
  \item Only 1 starch/sucrose/glycogen module: M00854 Glycogen biosynthesis (no dedicated trehalose-biosynthesis module in KEGG --- consistent with the paper's central complaint).
  \item vaa00500 confirmed: NAME=`Starch and sucrose metabolism - \textit{Variovorax} sp.\ PAMC 28711'.
\end{itemize}

\subsection{TreY coordinate audit}
The paper text reports ``\ldots started and stopped at 335612 to 3352054 coding sequence'' (missing one digit). Two clean digit restorations:
\begin{tabularx}{\linewidth}{@{}l l X@{}}
\toprule
Reading & Coordinates & Matches \\
\midrule
Missing `6': 3356112 to 3352054 & 3{,}352{,}054..3{,}356{,}112 & BV-BRC RAST peg.3325 exactly (4{,}059 nt, 1{,}352 aa) \\
Missing `7': 3357112 to 3352054 & 3{,}352{,}054..3{,}357{,}112 & PGAP AX767\_16200 within 7 bp (3{,}352{,}054..3{,}357{,}119) \\
\bottomrule
\end{tabularx}
Either reading identifies the same locus; the discrepancy is a typesetting error in the paper, not an analysis error.

\section{Coverage and Agreement Scoring}
\begin{tabularx}{\linewidth}{@{}l l X@{}}
\toprule
Category & Count & Notes \\
\midrule
Total testable claims enumerated       & 37 & up from 16 in pass-1 \\
Inherently untestable on free compute  & 11 & 5 MetaCyc + 6 historical DB snapshots \\
Testable on free compute               & 26 & score denominator \\
Tested in re-pass                      & 25 & only Han 2016 opine claim unverified \\
Verified                               & 21 & KEGG / RAST / PGAP / BV-BRC / TreX / glycogen / isolation / vaa00500 \\
Partial                                & 4  & claims 25 (typo), 27 (TreY pseudogene), 33 (pathway V incomplete), 8 (TreY MetaCyc inferred) \\
Contradicted                           & 0  & none \\
\bottomrule
\end{tabularx}

\medskip
\noindent\textbf{Re-pass scores (out of 10)}:
\begin{itemize}[leftmargin=1.2em,itemsep=1pt]
  \item \textbf{COVERAGE = 9/10} (up from 6/10 in pass-1). Lift: 21 additional enumerated claims, TreX added, glycogen cluster surveyed, genome features re-grounded from GenBank, Antarctic/lichen isolation source verified, vaa00500 + KEGG stats grounded by direct \texttt{rest.kegg.jp/get}. Held back by (i) MetaCyc permanently blocked, (ii) 2018 historical DB snapshots not queryable from current APIs.
  \item \textbf{AGREEMENT = 9/10} (from pass-1 8/10). Only ``partial'' items are a known paper typo (25) and the TreY pseudogene biological caveat (27, 33) the paper itself glosses.
\end{itemize}

\section{Verdict --- PARTIAL}
The core methodological claim (annotation-system divergence at TreY between KEGG/MetaCyc and RAST) fully replicates and is additionally corroborated by PGAP's pseudogene flag on the same locus. The paper's broader biological framing (``three functional trehalose-biosynthesis pathways in PAMC28711'') is not fully supported by the re-pass, because TreY is annotated as a frameshifted pseudogene by PGAP and returns no vaa hit for KEGG K06044. MetaCyc rows and 2018-snapshot statistics remain out of reach on free compute.

\section*{GENUINE CRITIQUE}
\addcontentsline{toc}{section}{GENUINE CRITIQUE}
\subsection*{What actually replicated}
\begin{itemize}[leftmargin=1.2em,itemsep=1pt]
  \item \textbf{Annotation-divergence thesis (the paper's real contribution)}: cleanly reproduces. KEGG \texttt{vaa} returns no gene for K06044 (TreY); RAST's SEED annotation on the same BV-BRC genome does return TreY at the expected coordinates. This is the paper's central methodological point and it survives re-pass unchanged.
  \item \textbf{Genome features}: 4{,}316{,}152 bp, 65.97\% GC, 1 circular chromosome, 4{,}104 CDS, 6 rRNA, 46 tRNA all match after direct recomputation from the GenBank flatfile.
  \item \textbf{Enzyme-presence pattern for OtsA, OtsB, TreZ, TreS}: fully verified via both KEGG KO$\to$vaa link queries \emph{and} PGAP product-name matches. Coordinates recorded per locus.
  \item \textbf{Isolation context}: Antarctic origin, lichen \textit{Himantormia} host, 2015 collection date, GCA\_001577265.1 assembly all verified against BV-BRC and PGAP qualifiers --- previously trusted without independent grounding.
  \item \textbf{vaa00500 pathway map} identity and the KEGG statement that no dedicated trehalose-biosynthesis \emph{module} exists in KEGG (only M00854 glycogen biosynthesis) both hold up under direct API queries.
\end{itemize}

\subsection*{What did NOT fully replicate}
\begin{itemize}[leftmargin=1.2em,itemsep=1pt]
  \item \textbf{``Three functional biosynthesis pathways''}: partial only. The TreY/TreZ pathway has functional TreZ and functional TreX and the full upstream glycogen cluster, but its central enzyme TreY (AX767\_16200) is flagged \texttt{/pseudo} by PGAP with a frameshift, and KEGG K06044 has zero vaa hits. The paper describes TreY as ``present in RAST'' (which is correct) but does not acknowledge PGAP's pseudogene call, which materially weakens the ``functional pathway'' framing.
  \item \textbf{MetaCyc column of Table 1}: not independently replicable on free compute. Pathway Tools is license-gated, no PAMC28711 PGDB exists publicly, and the paper's MetaCyc procedure is under-described (no PGDB build steps, no version beyond ``v22.5 Aug 2018''). Five of the fifteen Table-1 cells therefore remain \texttt{NOT\_TESTED}, not verified.
  \item \textbf{Table 2 historical statistics} (2{,}688 pathways / 339 modules / 381 superpathways / 530 maps / 15{,}329 vs 11{,}004 reactions): all six are August-2018 snapshots and cannot be re-derived from current APIs. This is not a fault of the paper --- database statistics necessarily drift --- but it caps re-pass coverage at 26/37.
  \item \textbf{TreY paper coordinates ``335612 to 3352054''}: partial. The coordinate string is missing one digit; two clean digit restorations both identify the same locus, so the biology is safe but the paper contains an uncorrected typesetting error.
  \item \textbf{Han 2016 opine-utilizing citation}: not verified. Out-of-scope external citation.
\end{itemize}

\subsection*{Methodological concerns with the paper}
\begin{itemize}[leftmargin=1.2em,itemsep=1pt]
  \item \textbf{No functional validation}: the paper's conclusions rest entirely on database annotation calls. No RT-qPCR, no proteomics, no growth assay under osmotic stress. The TreY pseudogene call would have been caught by any in-vitro validation and this weakens the ``pathway present'' language.
  \item \textbf{TreX omitted from the comparison table}: MetaCyc pathway V explicitly requires TreX (EC 3.2.1.68). The paper's Fig 2A names it but Table 1 does not include a TreX row, so the tri-database comparison is incomplete for pathway V. Re-pass added TreX and confirmed two functional copies in vaa.
  \item \textbf{MetaCyc procedure under-described}: no PGDB build details, no lineage traceability. Anyone attempting exact replication of the MetaCyc column would need to obtain a Pathway Tools license and rebuild a PGDB from the 2018 database snapshot --- practically impossible.
  \item \textbf{Uncorrected typo in TreY coordinates}: passed peer review despite one-digit dropout.
  \item \textbf{No acknowledgement of the TreY frameshift}: PGAP's \texttt{/pseudo} qualifier on AX767\_16200 is publicly available and should have been discussed alongside the ``present in RAST'' finding.
\end{itemize}

\subsection*{Overall assessment}
The paper's methodological contribution (documenting annotation-system divergence for a real biological example) is sound and replicates. Its biological framing overstates the ``three functional pathways'' conclusion by not addressing the pseudogene status of TreY. On free compute, MetaCyc-column re-verification is structurally blocked, so a complete numeric replication is not achievable regardless of effort. Confidence in the paper's core thesis is high; confidence in its biological interpretation is moderate.

\section{Artifacts}
See the Markdown source \texttt{report/REPORT.md} \S 7 for the full artifact inventory. Key JSON:
\texttt{results/repass/genome\_features.json}, \texttt{trex\_and\_cluster.json},
\texttt{kegg\_crosscheck.json}, \texttt{bvbrc\_metadata.json},
\texttt{claims\_enumerated.json}, \texttt{parser\_provenance.json}.

\vspace{1em}
\hrule
\vspace{0.5em}
\noindent\textit{Re-pass generated 2026-06-23 (CDT); LaTeX rendering 2026-07-05.
Tools: BioPython, KEGG REST API, BV-BRC API, poppler pdftotext, all on free compute.
No LLM-generated numbers --- every quantitative claim above is grounded in a JSON artefact under}
\texttt{results/repass/}.

\end{document}
